\subsection{Operational Envelope: Grid-Scale Analysis}
\label{sec:scalability}

A multi-condition sweep evaluated all three agents across four representative grid scales
($N = 20$ sensors, 20 seeds per condition) to characterise the operational envelope.
Table~\ref{tab:grid_summary} reports cumulative reward, Jain's fairness index, and energy
efficiency across grid sizes.

\begin{table}[H]
\centering
\caption{Mean cumulative reward ($\times 10^6$), Jain's fairness index $\mathcal{J}$, and
         energy efficiency (B/Wh), $N = 20$ sensors, 20 seeds per condition.
         All agents achieve 100\% sensor coverage at every grid scale.}
\label{tab:grid_summary}
\begin{tabularx}{\textwidth}{l l Z Z Z Z Z Z}
\toprule
Grid & Agent & Reward ($\times10^6$) & $\pm$Std & Mean $\mathcal{J}$ & $\pm$Std & Eff (B/Wh) & $\pm$Std \\
\midrule
\multirow{3}{*}{$100{\times}100$}
 & DQN             & 6.48 & 0.03 & 0.999 & 0.001 & 219.3 & 2.1 \\
 & SF-Aware Greedy & 6.50 & 0.02 & 0.999 & 0.001 & 220.0 & 1.8 \\
 & Nearest Greedy  & 6.49 & 0.03 & 0.998 & 0.001 & 219.6 & 2.0 \\
\midrule
\multirow{3}{*}{$300{\times}300$}
 & DQN             & 6.58 & 0.07 & 0.999 & 0.001 & 221.5 & 3.1 \\
 & SF-Aware Greedy & 6.56 & 0.07 & 0.999 & 0.001 & 220.9 & 3.0 \\
 & Nearest Greedy  & 6.57 & 0.07 & 0.999 & 0.001 & 221.3 & 2.9 \\
\midrule
\multirow{3}{*}{$500{\times}500$}
 & DQN             & 5.56 & 0.40 & 0.961 & 0.027 & 187.9 & 18.1 \\
 & SF-Aware Greedy & 5.60 & 0.39 & 0.964 & 0.026 & 189.1 & 17.6 \\
 & Nearest Greedy  & 5.57 & 0.41 & 0.960 & 0.028 & 188.1 & 18.4 \\
\midrule
\multirow{3}{*}{$1{,}000{\times}1{,}000$}
 & DQN             & 4.62 & 0.41 & 0.879 & 0.038 & 155.9 & 13.8 \\
 & SF-Aware Greedy & 4.79 & 0.39 & 0.885 & 0.036 & 161.2 & 13.1 \\
 & Nearest Greedy  & 4.62 & 0.42 & 0.876 & 0.039 & 155.8 & 14.0 \\
\bottomrule
\end{tabularx}
\end{table}

The operational envelope analysis presents a consistent and strong result across all grid scales:
all three agents achieve complete sensor coverage ($\text{coverage} = 100\%$) at every grid size
evaluated. This finding confirms that the domain-randomised DQN policy generalises effectively
across the full range of tested environments without retraining.

\paragraph{Trivial and near-trivial regimes (100$\times$100 and 300$\times$300).}
At these scales, all agents produce near-identical outcomes. On the $100 \times 100$ grid,
cumulative rewards cluster tightly around $6.49 \times 10^6$ and Jain's indices approach
$\mathcal{J} = 0.999$ for all agents, differences amounting to less than 0.3\%. On the
$300 \times 300$ grid, performance remains equally indistinguishable ($\mathcal{J} \approx 0.999$,
efficiency $\approx 221$\,B/Wh across all three). At these scales the energy budget is sufficient
for any routing strategy to visit all sensors and accumulate near-maximal data, leaving no room
for algorithmic differentiation.

\paragraph{Intermediate regime (500$\times$500).}
Performance diverges slightly. SF-Aware Greedy achieves a marginally higher Jain's fairness
index ($\mathcal{J} = 0.964$ vs.\ DQN $0.961$) and energy efficiency
($189.1$ vs.\ $187.9$\,B/Wh), a difference of approximately 1.6\%. The DQN and Nearest Greedy
are statistically comparable at this scale. The divergence is modest and within the inter-seed
variance, but reflects the onset of resource contention as the UAV must allocate a growing
proportion of flight steps to transit between sensors.

\paragraph{Large-scale regime (1,000$\times$1,000).}
At the largest grid scale, SF-Aware Greedy achieves a more pronounced efficiency advantage
($161.2$ vs.\ $155.9$\,B/Wh for DQN, a difference of 3.4\%) and marginally higher fairness
($\mathcal{J} = 0.885$ vs.\ $0.879$). The DQN and Nearest Greedy are near-identical at this
scale. The SF-Aware Greedy's advantage arises from its explicit criterion of routing toward
the lowest spreading-factor sensor, which minimises per-visit flight cost at large inter-sensor
distances. Despite the performance gap, all agents maintain 100\% coverage, confirming that
physical feasibility is not violated at this scale. The domain-randomised DQN, requiring no
hand-crafted routing heuristic, achieves performance within 3.5\% of the best baseline across
the entire grid-size range.
